% front.tex (국문 v13): 제목 · 저자 · 초록 · 핵심어(블록 1), 자료 가용성(블록 3), 뒷부분(블록 4).
\providecommand{\frontblock}{0}
\ifnum\frontblock=1
\title[물리 잔차 기반 기계학습의 연도별 활동층 두께 예측]{물리 잔차 기반 기계학습 워크플로를 활용한 영구동토 활동층 두께 예측과 지역별 편차 분석}

\author*[1]{\fnm{백승원}}\email{[교신 저자 전자우편]}
\author[2]{\fnm{[공저자]} \sur{[극지연구소]}}
\affil*[1]{\orgdiv{에너지시스템공학부}, \orgname{서울대학교}, \orgaddress{\city{서울}, \country{대한민국}}}
\affil[2]{\orgname{극지연구소}, \orgaddress{\city{인천}, \country{대한민국}}}

\abstract{활동층 두께(ALT)는 해마다 달라지지만, 기존 격자 ALT 제품은 대부분 여러 해의 조건으로 학습되었고 같은 위치 · 같은 해의 관측과 비교된 적이 드물다. 이 연구는 알래스카, 레나 델타, 캐나다 북서부의 지점-연도 ALT 관측 \nncyAll{}개(2010--2024년)에 관측일까지 누적한 ERA5-Land 융해 도일을 결합하고, 지역 라벨로 재보정한 Stefan 모델의 잔차를 gradient boosting 으로 학습하는 물리 잔차 워크플로를 검증하였다. 공간 블록 교차검증에서 물리 잔차 모델은 Stefan 모델의 오차를 레나 델타에서 \ndrefLeblkRoabs{} cm, 캐나다에서 \ndrefCablkRoabs{} cm 줄였고 알래스카에서는 같은 수준이었다. 레나 델타를 학습에서 제외하면 직접 기계학습의 오차는 Stefan 모델보다 \ndrefLeregDzabs{} cm 컸으나, 물리 잔차 모델은 Stefan 모델 수준을 유지하였다. 지역 라벨 25개를 더하면 물리 잔차 모델의 레나 델타 오차는 \ntlrmseLeRszero{} cm에서 \ntlrmseLeRstwfive{} cm로 줄었고, 외부 CALM 지점에서는 라벨 5개로 오차가 \nftAllRszero{} cm에서 \nftAllRsfive{} cm로 감소하였다. 같은 셀 · 같은 해 비교에서 레나 델타의 오차(\nccLeRoall{} cm)는 공개 제품 5종(\nccLeCcifiveall--\nccLeWeiall{} cm)보다 낮았다. 따라서 Stefan 모델을 기준으로 두고 잔차만 학습하는 구조는 대상 지역 라벨이 적거나 없을 때 가장 유용하다. 이 워크플로는 적용 가능 영역이 명시된 2011--2024년의 연도별 1 km ALT 지도를 제공한다.}

\keywords{영구동토, 활동층 두께, Stefan 모델, 물리 기반 기계학습, 공간 교차검증, ERA5-Land}
\fi

\ifnum\frontblock=3
\bmhead{자료 가용성}
모든 입력 자료는 공개 자료이다. ALT 관측: ABoVE 토양 수분 · 활동층 두께 v2(ORNL DAAC, \url{https://doi.org/10.3334/ORNLDAAC/2369}), ALLena 융해 깊이 자료(PANGAEA, \url{https://doi.org/10.1594/PANGAEA.973813}), GTN-P CALM(PANGAEA, \url{https://doi.org/10.1594/PANGAEA.972777}). 기후: ERA5-Land 월평균(Copernicus Climate Data Store). 공변량: Copernicus DEM GLO-30, SoilGrids 2.0, ESA CCI 영구동토 확률, MODIS MOD13Q1 · MOD10A2 v6.1(AppEEARS). 제품: Wei 등(\url{https://doi.org/10.5281/zenodo.21667583}), Liu 등(\url{https://doi.org/10.5281/zenodo.10070610}), ESA CCI ALT v4 · v5(CEDA), Aalto 등(\url{https://zenodo.org/records/5003804}), Yi와 Kimball(\url{https://doi.org/10.3334/ORNLDAAC/1760}). 지점-연도 라벨 표, 연도별 기후 입력 표, 정적 공변량 표, 모든 방법의 겹 밖 예측, 알래스카 · 레나 델타의 연도별 1 km 지도(2011--2024년)는 게재 시 공개 저장소에 둔다[저장소 · DOI].
\fi

\ifnum\frontblock=4
\bmhead{감사의 글}
[연구비 · 감사의 글.]

\bmhead{저자 기여}
백승원은 연구를 설계하고 자료를 구축하였으며 분석을 수행하고 원고를 작성하였다. [공저자 기여.]

\bmhead{추가 정보}
\textbf{보충 자료} 보충 방법 S1--S5, 보충 표 S1--S6, 보충 그림 S1--S10을 함께 제공한다.

\textbf{이해 상충} 저자들은 이해 상충이 없음을 밝힌다.
\fi
